\section{Methodology}
\label{sec:methodology}

Our investigation follows a comparative empirical approach to test the hypothesis that attention is a fundamental bottleneck mechanism across both AI and Internet systems. We evaluate this across two representative tasks: sequential recommendation and click-through rate (CTR) prediction.

\subsection{Sequential Recommendation (MovieLens)}
The goal of sequential recommendation is to predict the next movie a user will interact with, given their history of interactions. We treat this as a sequence prediction task.

\para{Dataset and Preprocessing}
We use the \textsc{MovieLens} 1M dataset, which contains approximately 1 million ratings from 6,000 users. For our experiments, we used a subset of 891,382 ratings from 11,923 users (after filtering for interaction sequences of at least 5 items). Interaction sequences are ordered by row index, representing a temporal stream of user-movie ratings.

\para{Models}
\begin{itemize}[leftmargin=*,itemsep=0pt,topsep=0pt]
    \item \textsc{GRUModel}: A standard Gated Recurrent Unit (GRU) as a representative non-attention sequential baseline.
    \item \textsc{TransformerModel}: A multi-head self-attention model using positional embeddings, representing the attention-based approach to sequence modeling.
\end{itemize}

\subsection{CTR Prediction (Criteo)}
The goal of CTR prediction is to estimate the probability that a user will click on an online advertisement, given both dense (numerical) and categorical features.

\para{Dataset and Preprocessing}
We use a 20,000-row subset of the \textsc{Criteo} CTR dataset. Categorical features are encoded using a label encoder, and numerical features are normalized using a MinMaxScaler. Missing values are filled with default values (0 for dense, -1 for categorical).

\para{Models}
\begin{itemize}[leftmargin=*,itemsep=0pt,topsep=0pt]
    \item \textsc{MLP}: A standard multi-layer perceptron (MLP) where all features (embedded categorical and normalized dense) are concatenated as input.
    \item \textsc{AttentionMLP}: An attention-based architecture that applies self-attention over feature embeddings to capture feature-level interactions before passing them through an MLP.
\end{itemize}

\subsection{Experimental Setup}
All models were implemented in PyTorch and trained using the Adam optimizer with a learning rate of 0.001. We use Hit Rate @10 (HR@10) for the MovieLens sequential recommendation task and validation accuracy for the Criteo CTR task. To ensure robustness, results are averaged over multiple epochs (3 for MovieLens, 5 for Criteo).

\subsection{Mechanism Analysis: Attention Entropy}
To examine the underlying mechanism, we compute the attention entropy $H(A)$ for each prediction:
\begin{equation}
    H(A) = - \sum_{i} a_i \log a_i
\end{equation}
where $a_i$ is the attention weight for the $i$-th element. We then analyze the relationship between attention entropy and prediction correctness to test whether "focused" attention (lower entropy) correlates with higher predictive accuracy.
